\section*{Summary}

An exact remaining gap is the strongest statement the book derives about a quantity, with its conditions, together with the single step that would turn it into the full result (\cref{sec:22.2}). The book contains 74. Three limit only a proof and leave every number standing; of the 71 that involve a number, 10 come with a bound or a bracket, 17 with an estimate, 39 are unsized and 5 are independent checks of a computed number. The missing step is a specified calculation in 56 of them, a measurement in 11 and a proof or method in 7. Several gaps share one missing step, a finite-momentum Bethe--Salpeter calculation for silicon and germanium or quasiparticle bands of germanium with spin--orbit coupling among them, so the 74 rows stand for fewer independent calculations; \cref{sec:22.1} orders those that would change a printed result, limit or ranking, each with the exact relation and the control against data that it must pass (\cref{tab:22-calc}). Unsized gaps dominate both the chapters from magnetic targets to the comparison across channels (16 of 19) and those on the measurement chain (14 of 17), where the rates of capture, trapping and release at millikelvin are neither measured nor computed.

The master formula holds wherever the dark-matter side can be traced out before the target (\cref{sec:22.3}): the dark field acts through species-conserving one-body rows, whatever leaves on the dark side is a free plane wave, and the target starts in a Gibbs state and is read through the excitation it is left in. Charged-current absorption, $\chi + (A, Z) \to (A, Z+1) + e^-$, breaks the first two. Its nuclear rows change the element, and its outgoing electron moves in the Coulomb field of the final state of the target, so the dark-matter side becomes an operator on the target. In the allowed approximation the rate is the relativistic Fermi function of the daughter charge times the phase space $p_{\mathrm{ej}}E_e$ times the Fermi and Gamow--Teller strengths. For $^{125}$Te at $\mchi = \SI{1}{MeV}$ the Fermi function is 6.34, \SI{4.5}{\percent} above its value with the charge of the parent, and it keeps the rate finite at threshold. The whole Fermi strength, $A - 2Z$, sits in the isobaric analog state 13 to \SI{14}{MeV} up, so Gamow--Teller transitions with a quenched axial coupling carry the rate below it, and a count of one conversion per neutron overstates the strength. The radiative decay of an excited dark state has no target operator in its amplitude; the track of a recoiling heavy ion is a transport problem whose stopping alone is a response; strong coupling and targets with coherences fail the assumptions of the formula. The collective modes of a charge-density wave lie inside the framework and need a conserving response whose proper part contains the Goldstone phason.

Most of what the book leaves open is calculation rather than theory. In 20 of its 22 channels (\cref{sec:22.1}) the response has a controlled form, the next term is known and its check is stated; what is missing there is the number, and a stated calculation or measurement supplies it.
